\documentclass{article}

\usepackage[preprint]{neurips_2026}

\usepackage[utf8]{inputenc}
\usepackage[T1]{fontenc}
\usepackage[hidelinks]{hyperref}
\usepackage{url}
\usepackage{booktabs}
\usepackage{amsfonts}
\usepackage{nicefrac}
\usepackage{microtype}
\usepackage[dvipsnames]{xcolor}
\usepackage{graphicx}
\usepackage{amsmath}
\usepackage{cleveref}
\usepackage{amssymb}
\usepackage{mathtools}
\usepackage{amsthm}
\usepackage{enumitem}
\usepackage{algorithm}
\usepackage{algorithmic}

\theoremstyle{plain}
\newtheorem{theorem}{Theorem}[section]
\newtheorem{proposition}[theorem]{Proposition}
\newtheorem{lemma}[theorem]{Lemma}
\newtheorem{corollary}[theorem]{Corollary}
\theoremstyle{definition}
\newtheorem{definition}[theorem]{Definition}
\newtheorem{assumption}[theorem]{Assumption}
\theoremstyle{remark}
\newtheorem{remark}[theorem]{Remark}

\newcommand{\SO}{\mathrm{SO}}
\newcommand{\SH}{Y}
\newcommand{\CG}[6]{C^{#5,\,#6}_{#1,\,#2;\,#3,\,#4}}
\newcommand{\abs}[1]{\lvert #1 \rvert}
\newcommand{\norm}[1]{\lVert #1 \rVert}
\newcommand{\R}{\mathbb{R}}
\newcommand{\C}{\mathbb{C}}
\newcommand{\conj}[1]{\overline{#1}}
\newcommand{\ip}[2]{\langle #1,\, #2 \rangle}
\DeclareMathOperator{\sgn}{sgn}

\title{The Selective $\SO(3)$ Bispectrum on $S^2$:\\
Construction and Completeness}

\author{%
  Johan Mathe \\
  Atmo \\
  San Francisco, CA \\
  \texttt{johan@atmo.ai}
}

\begin{document}

\maketitle

\begin{abstract}
We construct a \emph{selective} bispectrum for $\SO(3)$ acting on the
$2$-sphere $S^2$ at band-limit $L$, using $\mathcal{O}(L^2)$ scalar
entries instead of the full $\mathcal{O}(L^3)$.
The construction recovers spherical harmonic coefficients degree by
degree: explicit gauge fixing resolves the continuous $\SO(3)$ ambiguity
at low degrees, after which a linear bootstrap---whose generic
full-rank property follows from a polynomial non-vanishing
argument---handles all higher degrees.
We prove that the resulting selective set is a complete invariant for
generic band-limited signals: two signals share the same selective
bispectrum if and only if they differ by an $\SO(3)$ rotation, for
Lebesgue-almost-every signal.
On an equiangular grid with $P$ pixels, the selective bispectrum
has $\mathcal{O}(P)$ entries and costs $\mathcal{O}(P^{3/2})$ to
compute, compared with $\mathcal{O}(P^{3/2})$ entries and
$\mathcal{O}(P^2)$ cost for the full bispectrum.
This extends the selective $G$-bispectrum of Mataigne et
al.~\cite{mataigne2024} from finite groups to the compact continuous
group $\SO(3)$.
\end{abstract}

\section{Introduction}\label{sec:introduction}

Mataigne et al.~\cite{mataigne2024} showed that for finite groups $G$,
the full $G$-bispectrum---which has $\mathcal{O}(|G|^2)$ matrix-valued
entries---can be reduced to a \emph{selective} subset of only
$\mathcal{O}(|G|)$ entries while preserving completeness. Their
construction traverses the irrep graph via a generating set, recovering
Fourier coefficients $\mathcal{F}(\Theta)_\rho$ one irrep at a time.

We extend this idea from finite groups to $\SO(3)$ acting on the
$2$-sphere $S^2$, where signals are expanded in spherical harmonics
(which form bases for the irreducible representations of $\SO(3)$) and bispectral entries are \emph{scalars}
rather than matrices. The passage from finite to continuous introduces
two key differences:

\begin{enumerate}[nosep]
  \item \textbf{Scalars vs.\ matrices.}
    On a finite group, $\mathcal{F}(\Theta)_\rho$ is a
    $d_\rho \times d_\rho$ matrix, and a single bispectral entry
    $\beta_{\rho_1,\rho_2}$ is itself matrix-valued, carrying
    $\mathcal{O}(d_\rho^2)$ scalar constraints. On $S^2$, the coefficient
    vector $\mathbf{F}_\ell$ at degree~$\ell$ has $2\ell+1$ components
    but each bispectral entry $\beta_{\ell_1,\ell_2,\ell}$ is a
    \emph{single scalar}. The chain
    $\beta_{1,0,1},\,\beta_{1,1,2},\,\dotsc$ provides only $\sim L$
    scalars---far too few to pin down all $(L+1)^2$ unknowns.

  \item \textbf{Continuous gauge.}
    On a finite group, the inversion ambiguity is a discrete element
    $h \in G$, absorbed by the CG structure. On $S^2$, the ambiguity
    is a continuous $\SO(3)$ rotation (3~real parameters), which we
    resolve by explicit gauge fixing: aligning $\mathbf{F}_1$ with the
    $z$-axis and fixing the azimuthal phase of~$\mathbf{F}_2$.
\end{enumerate}

\paragraph{High-level construction.}
Our selective bispectrum recovers the spherical harmonic coefficients
$\mathbf{F}_0, \mathbf{F}_1, \dotsc, \mathbf{F}_L$ degree by degree,
analogous to Mataigne et al.'s irrep-by-irrep traversal but adapted
for the scalar bispectrum on~$S^2$.  At each degree~$\ell$, we need
$2\ell+1$ scalar equations to determine $\mathbf{F}_\ell$. We obtain
these from bispectral triples in which $\mathbf{F}_\ell$ appears
linearly (the other two indices being lower, hence already known). The
key insight is:

\begin{quote}
  \emph{For $\ell \ge 4$, the number of ``chain'' and ``cross'' triples
  involving degree~$\ell$ linearly grows as $\Theta(\ell^2)$, which
  exceeds the $2\ell+1$ equations needed. The entire recovery from
  degree~$4$ onwards is therefore a linear solve. Only the base cases
  $\ell = 0, 1, 2, 3$ require nonlinear equations and gauge fixing
  (with a remaining open question at $\ell = 2$; see
  \Cref{rem:ell2-gap}).}
\end{quote}

\noindent
The total number of selected triples sums to $(L+1)^2 - 1
= \mathcal{O}(L^2)$, matching the dimension of the orbit space
(hence optimal). In pixel space ($P \approx 4L^2$ samples on $S^2$),
this is $\mathcal{O}(P)$---the selective bispectrum is linear in the
input size.


\section{Setup and Notation}\label{sec:setup}

\begin{definition}[Band-limited signal on $S^2$]
A signal $f : S^2 \to \R$ is \emph{band-limited} at degree $L$ if its
spherical harmonic expansion
\begin{equation}\label{eq:sh-expansion}
  f(\theta, \varphi)
  = \sum_{\ell=0}^{L} \sum_{m=-\ell}^{\ell}
    a_\ell^m \, \SH_\ell^m(\theta, \varphi)
\end{equation}
has $a_\ell^m = 0$ for all $\ell > L$. We write
$\mathbf{F}_\ell = (a_\ell^{-\ell}, \dotsc, a_\ell^{\ell})
  \in \C^{2\ell+1}$
for the coefficient vector at degree~$\ell$.
\end{definition}

\noindent
For real $f$, the conjugation symmetry
$a_\ell^{-m} = (-1)^m \conj{a_\ell^m}$ holds, so
$\mathbf{F}_\ell$ has $2\ell+1$ real degrees of freedom.
The total signal dimension is
$\dim = \sum_{\ell=0}^{L}(2\ell+1) = (L+1)^2$.

\begin{definition}[Rotation action]\label{def:rotation-action}
Under $g \in \SO(3)$, the spherical harmonic coefficients transform as
$\mathbf{F}_\ell \mapsto D^\ell(g)\,\mathbf{F}_\ell$,
where $D^\ell$ is the $(2\ell+1)$-dimensional Wigner~$D$-matrix.
\end{definition}

\begin{definition}[Bispectrum on $S^2$]\label{def:bispectrum}
For each triple $(\ell_1, \ell_2, \ell)$ satisfying the triangle
inequality $\abs{\ell_1 - \ell_2} \le \ell \le \ell_1 + \ell_2$,
the \emph{bispectral coefficient} is
\begin{equation}\label{eq:bispectrum}
  \beta_{\ell_1,\ell_2,\ell}
  = \sum_{\substack{m_1,\,m_2,\,m \\ m_1+m_2=m}}
    \CG{\ell_1}{m_1}{\ell_2}{m_2}{\ell}{m}\;
    a_{\ell_1}^{m_1}\,a_{\ell_2}^{m_2}\,\conj{a_\ell^m},
\end{equation}
where $\CG{\ell_1}{m_1}{\ell_2}{m_2}{\ell}{m}$ denotes the
Clebsch--Gordan coefficient
$\langle \ell_1, m_1;\, \ell_2, m_2 \mid \ell, m \rangle$.
Each $\beta_{\ell_1,\ell_2,\ell}$ is a scalar (real for real~$f$).
\end{definition}

\begin{proposition}[Size of the full bispectrum]\label{prop:full-size}
The number of admissible triples $(\ell_1, \ell_2, \ell)$ with
$0 \le \ell_1, \ell_2, \ell \le L$ satisfying the triangle inequality
$|\ell_1 - \ell_2| \le \ell \le \ell_1 + \ell_2$ is $\Theta(L^3)$.
\end{proposition}
\begin{proof}
\emph{Strategy: we count the admissible lattice points
$(\ell_1, \ell_2, \ell)$ by summing the number of valid~$\ell$ for each
pair $(\ell_1, \ell_2)$, then evaluate the resulting double sum in
closed form.}

For fixed $\ell_1$ and $\ell_2$, the triangle inequality admits
$\ell \in \{|\ell_1 - \ell_2|, \dotsc, \min(\ell_1 + \ell_2, L)\}$,
giving $\min(\ell_1 + \ell_2, L) - |\ell_1 - \ell_2| + 1$ valid
values of $\ell$. For $\ell_1, \ell_2 \le L/2$, this is
$2\min(\ell_1, \ell_2) + 1$. Summing:
\[
  N_{\mathrm{full}}
  = \sum_{\ell_1=0}^{L}\;\sum_{\ell_2=0}^{L}\;
    \bigl(\min(\ell_1+\ell_2,L) - |\ell_1-\ell_2| + 1\bigr)
  = \tfrac{1}{3}L^3 + \mathcal{O}(L^2).
\]
To see this, note that the dominant contribution comes from the
interior of the summation region where $\ell_1 + \ell_2 \le L$.
In that regime each pair $(\ell_1, \ell_2)$ contributes
$2\min(\ell_1, \ell_2) + 1$ triples. Summing over
$0 \le \ell_1 \le \ell_2$ and $\ell_1 + \ell_2 \le L$ (then doubling
by symmetry) gives
\[
  2\sum_{\ell_1=0}^{\lfloor L/2\rfloor}\;
    \sum_{\ell_2=\ell_1}^{L-\ell_1}(2\ell_1 + 1)
  = 2\sum_{\ell_1=0}^{\lfloor L/2\rfloor}
    (2\ell_1+1)(L - 2\ell_1 + 1)
  = \tfrac{1}{3}L^3 + \mathcal{O}(L^2),
\]
where the last equality follows from standard summation identities.
Each such triple yields exactly one scalar $\beta_{\ell_1,\ell_2,\ell}$,
establishing the $\Theta(L^3)$ count.
\end{proof}

\begin{remark}[Ordered vs.\ unordered triples]
For real-valued signals, $\beta_{\ell_1,\ell_2,\ell} =
\overline{\beta_{\ell_2,\ell_1,\ell}}$, so only unordered pairs
$\{\ell_1,\ell_2\}$ contribute independent entries, reducing the count
by a factor of~$2$ (to $\sim L^3/6$). The $\Theta(L^3)$ scaling is
unaffected, and for our selective construction we work with ordered
triples throughout, which simplifies the category analysis.
\end{remark}

\begin{remark}[Invariance]
Each $\beta_{\ell_1,\ell_2,\ell}$ is $\SO(3)$-invariant: the Wigner
matrices satisfy $D^{\ell_1} \otimes D^{\ell_2}
= \bigoplus_\ell D^\ell$ under the CG decomposition, so the triple
product contracted with CG coefficients is unchanged by
$\mathbf{F}_\ell \mapsto D^\ell(g)\mathbf{F}_\ell$.
\end{remark}

\begin{theorem}[Kakarala~\cite{kakarala2009}]\label{thm:full-complete}
The full bispectrum $\{\beta_{\ell_1,\ell_2,\ell}\}$ determines $f$
up to $\SO(3)$ rotation for generic band-limited~$f$.
\end{theorem}


\section{Why $\mathcal{O}(L)$ Is Impossible}\label{sec:lower-bound}

Before constructing the selective set, we establish that
$\mathcal{O}(L^2)$ is a hard lower bound---no reduction to
$\mathcal{O}(L)$ is possible.

For finite groups of order $|G|$, the selective bispectrum achieves
$\mathcal{O}(|G|)$ scalar entries~\cite{mataigne2024}. On $S^2$, one
might hope that the ``BFS chain'' $\beta_{1,0,1},\,
\beta_{1,1,2},\,\dotsc,\,\beta_{1,L-1,L}$---just $L$ scalar
entries---suffices.

\begin{proposition}[Dimensional lower bound]\label{prop:lower-bound}
Any complete $\SO(3)$-invariant for band-limited signals on $S^2$
at band-limit $L$ must have at least $(L+1)^2 - 3$ functionally
independent components.
\end{proposition}
\begin{proof}
\emph{Strategy: a dimension-counting argument on the orbit space.
A complete invariant must separate orbits, hence must have at least
as many independent components as the dimension of the quotient
$\R^{(L+1)^2}/\SO(3)$.}

The orbit space $\R^{(L+1)^2} / \SO(3)$ has dimension
$(L+1)^2 - \dim\SO(3) = (L+1)^2 - 3$. Any injective map from the
orbit space to $\R^k$ requires $k \ge (L+1)^2 - 3$, hence
$\Omega(L^2)$.
\end{proof}

\begin{remark}
The chain $\beta_{1,\ell,\ell+1}$ for $\ell = 0, \dotsc, L-1$ provides
only $\sim 3L$ scalar constraints (since each chain entry couples
$\mathbf{F}_1 \in \C^3$ with one known $\mathbf{F}_\ell$). This cannot
recover $(L+1)^2 \approx L^2$ unknowns.
\end{remark}


\section{The Selective Index Set}\label{sec:index-set}

The goal is to select, for each target degree $\ell$, a set of
bispectral triples that provides enough equations to recover
$\mathbf{F}_\ell$ from already-known lower-degree coefficients. Since
$\mathbf{F}_\ell$ has $2\ell+1$ real degrees of freedom, we need at
most $2\ell+1$ entries. We classify candidate triples by how
$\mathbf{F}_\ell$ enters the bispectral formula~\eqref{eq:bispectrum}:
\emph{linearly} (when it occupies a position where the other two
indices are known), or \emph{nonlinearly} (when it appears more than
once). We prioritise linear entries because they yield a well-posed
linear system; nonlinear entries are used only at low degrees where
the pool of linear candidates is too small.

\begin{definition}[Selective bispectrum on $S^2$]\label{def:selective}
For each target degree $\ell = 0, \dotsc, L$, we select at most
$2\ell+1$ bispectral triples from four categories, in the following
priority order.

\medskip
\noindent\textbf{Category~1 (Chain).}
Triples $(\ell_1, \ell_2, \ell)$ with
$\ell_1 \le \ell_2 < \ell$ and $\ell_1 + \ell_2 \ge \ell$.
These give equations \emph{linear} in $\mathbf{F}_\ell$: the unknown
$\conj{a_\ell^m}$ appears once, multiplied by known lower-degree
coefficients.

\medskip
\noindent\textbf{Category~2 (Cross).}
Triples $(\ell_1, \ell, \ell')$ with
$1 \le \ell_1 < \ell$ and $\ell - \ell_1 \le \ell' < \ell$.
Here $\mathbf{F}_\ell$ appears in the second (tensor-product) position,
yielding equations linear in $a_\ell^{m_2}$ with known $a_{\ell_1}^{m_1}$
and $a_{\ell'}^{m}$.

\medskip
\noindent\textbf{Category~3 (Power).}
The triple $(0, \ell, \ell)$, giving
$\beta_{0,\ell,\ell} = a_0^0 \norm{\mathbf{F}_\ell}^2$
(a \emph{quadratic} constraint). Selected only when the linear
categories provide fewer than $2\ell+1$ equations.

\medskip
\noindent\textbf{Category~4 (Self-coupling).}
Triples $(\ell, \ell, \ell')$ with $0 \le \ell' < \ell$.
These are quadratic in $\mathbf{F}_\ell$. Selected only for the remaining deficit at low
degrees.

\medskip
\noindent
At each degree~$\ell$, entries are drawn by priority until $2\ell+1$
entries are reached (or all candidates are exhausted).
\end{definition}

The precise selection procedure is given in
Algorithm~\ref{alg:selective}, which mirrors the implementation in
\texttt{SO3onS2(selective=True)}. Within each category, candidates are
enumerated in \emph{descending} degree order (highest $\ell_2$ first
for chain, highest $\ell_1$ first for cross), which empirically yields
better-conditioned linear systems at each recovery step.

\begin{algorithm}[t]
\caption{Selective index set $\mathcal{S}$ for the $\SO(3)$ bispectrum
  on $S^2$ at band-limit~$L$.}
\label{alg:selective}
\small
\begin{algorithmic}[1]
\STATE \textbf{Input:} Band-limit $L$
\STATE \textbf{Output:} Index set
  $\mathcal{S} \subset \{(\ell_1,\ell_2,\ell) :
  |\ell_1-\ell_2| \le \ell \le \ell_1+\ell_2,\;
  0 \le \ell_1,\ell_2,\ell \le L\}$
\STATE $\mathcal{S} \leftarrow \{(0,0,0)\}$
\FOR{$\ell = 1, \dotsc, L$}
  \STATE $\text{budget} \leftarrow 2\ell + 1$
  \STATE $\text{candidates} \leftarrow [\,]$
    \hfill\COMMENT{ordered list}
  \STATE
  \STATE \COMMENT{\textbf{Category 1: Chain}
    --- linear in $\mathbf{F}_\ell$ (via $\conj{a_\ell^m}$)}
  \FOR{$\ell_2 = \ell{-}1, \dotsc, 0$}
    \FOR{$\ell_1 = \min(\ell_2,\, \ell{-}1), \dotsc, 0$}
      \IF{$\ell_1 + \ell_2 \ge \ell$}
        \STATE append $(\ell_1,\, \ell_2,\, \ell)$ to candidates
      \ENDIF
    \ENDFOR
  \ENDFOR
  \STATE
  \STATE \COMMENT{\textbf{Category 2: Cross}
    --- linear in $\mathbf{F}_\ell$ (via $a_\ell^{m_2}$)}
  \FOR{$\ell_1 = \ell{-}1, \dotsc, 1$}
    \FOR{$\ell' = \ell{-}1, \dotsc, \ell{-}\ell_1$}
      \STATE append $(\ell_1,\, \ell,\, \ell')$ to candidates
          \hfill\COMMENT{$\ell_1{+}\ell > \ell{-}1$ always}
    \ENDFOR
  \ENDFOR
  \STATE
  \STATE \COMMENT{\textbf{Category 3: Power}
    --- quadratic ($a_0^0\,\|\mathbf{F}_\ell\|^2$)}
  \STATE append $(0,\, \ell,\, \ell)$ to candidates
  \STATE
  \STATE \COMMENT{\textbf{Category 4: Self-coupling}
    --- quadratic/cubic}
  \FOR{$\ell' = \ell{-}1, \dotsc, 0$}
    \STATE append $(\ell,\, \ell,\, \ell')$ to candidates
  \ENDFOR
  \STATE
  \STATE Deduplicate candidates (preserving order)
  \STATE $\mathcal{S} \leftarrow \mathcal{S}
    \;\cup\;$ first $\min(\text{budget},\,
    |\text{candidates}|)$ entries
\ENDFOR
\RETURN $\mathcal{S}$
\end{algorithmic}
\end{algorithm}

\begin{proposition}[Entry counts]\label{prop:entry-counts}
Let $N_{\mathrm{chain}}(\ell)$ and $N_{\mathrm{cross}}(\ell)$ denote
the number of chain and cross candidates at degree~$\ell$.
\begin{enumerate}[label=(\roman*),nosep]
  \item For $\ell \ge 4$,
    $N_{\mathrm{chain}}(\ell) + N_{\mathrm{cross}}(\ell) \ge 2\ell + 1$.
    Only linear entries are selected.
  \item For $\ell \le 3$, additional power and self-coupling entries are
    required to reach $2\ell+1$.
\end{enumerate}
The per-degree counts are:
\begin{center}
\small
\begin{tabular}{cccccc}
  \toprule
  $\ell$ & Chain & Cross & Power & Self & Total (cap $2\ell+1$)\\
  \midrule
  0 & 0 & 0 & 0 & 1 & 1 \\
  1 & 0 & 0 & 1 & 1 & 2 \\
  2 & 1 & 1 & 1 & 2 & 5 \\
  3 & 2 & 3 & 1 & 1 & 7 \\
  $\ge 4$ & \multicolumn{2}{c}{$\ge 2\ell+1$} & 0 & 0 & $2\ell+1$ \\
  \bottomrule
\end{tabular}
\end{center}
\end{proposition}
\begin{proof}
\emph{Strategy: direct enumeration of lattice points.
We count the integer pairs $(\ell_1, \ell_2)$ satisfying each
category's constraints, verify the totals for small~$\ell$, and show
that the quadratic growth of both chain and cross counts exceeds
$2\ell+1$ for $\ell \ge 4$.}

Chain candidates at degree~$\ell$: pairs
$(\ell_1, \ell_2)$ with $\ell_1 \le \ell_2 < \ell$ and
$\ell_1 + \ell_2 \ge \ell$. Cross candidates: pairs
$(\ell_1, \ell')$ with $1 \le \ell_1 < \ell$ and
$\ell - \ell_1 \le \ell' < \ell$.

For $\ell = 4$: $N_{\mathrm{chain}} = 4$,
$N_{\mathrm{cross}} = 6$, total $= 10 > 9 = 2\ell+1$.
For $\ell = 5$: $N_{\mathrm{chain}} = 6$,
$N_{\mathrm{cross}} = 10$, total $= 16 > 11$.
In general, $N_{\mathrm{chain}}(\ell) = \Theta(\ell^2)$ (the number
of pairs $(\ell_1,\ell_2)$ with $\ell_1 \le \ell_2 < \ell$ and
$\ell_1 + \ell_2 \ge \ell$)
and $N_{\mathrm{cross}}(\ell) = \binom{\ell}{2} =
\frac{\ell(\ell-1)}{2}$; both grow as $\Theta(\ell^2)$,
exceeding $2\ell+1$ for $\ell \ge 4$.
\end{proof}

\begin{proposition}[Output size]\label{prop:output-size}
The selective bispectrum has
$\displaystyle
  \sum_{\ell=0}^{L} \min\bigl(2\ell+1,\;
    \text{entries at } \ell\bigr)
  = (L+1)^2 - 1
$
scalar entries. This matches the dimensional lower bound of
\Cref{prop:lower-bound} up to an additive constant.
\end{proposition}
\begin{proof}
\emph{Strategy: direct summation, splitting the low-degree base cases
($\ell \le 3$) from the $2\ell+1$ regime ($\ell \ge 4$) and using
the identity $\sum_{\ell=0}^{L}(2\ell+1) = (L+1)^2$.}

$1 + 2 + 5 + 7 + \sum_{\ell=4}^{L}(2\ell+1)
= 15 + \bigl[(L+1)^2 - 16\bigr] = (L+1)^2 - 1$.
\end{proof}


\section{Completeness}\label{sec:completeness}

The proof follows the same ``frequency-marching'' philosophy as
Mataigne et al.'s inversion algorithm for finite groups
(\cite{mataigne2024}, Algorithms~1--2): recover the Fourier
coefficients one irrep (here: one degree~$\ell$) at a time, using
bispectral entries that couple the unknown degree with already-known
ones.  The main structural difference is that on $S^2$ each bispectral
entry is a scalar rather than a matrix, so we need $2\ell+1$ entries
per degree instead of one matrix entry per irrep.  We state the
genericity conditions, the main theorem, and then establish each
recovery step as a separate lemma.

\begin{assumption}\label{ass:generic}
The signal $f$ satisfies:
\begin{enumerate}[label=(G\arabic*),nosep]
  \item $a_0^0 \ne 0$ (non-zero mean).
  \item $\norm{\mathbf{F}_1} \ne 0$ ($f$ has a non-zero dipole component).
  \item After gauge-fixing (\Cref{lem:gauge}), $a_2^1 \ne 0$ (azimuthal
    gauge is resolvable).
  \item For each $\ell \ge 4$, the linear system of
    \Cref{lem:linear-bootstrap} has full rank.
\end{enumerate}
\end{assumption}

\begin{remark}
Each condition excludes a proper algebraic subvariety of the signal
space $\R^{(L+1)^2}$. Their union is therefore Lebesgue-null.
\end{remark}

\begin{theorem}[Completeness of the selective $\SO(3)$ bispectrum on $S^2$]
\label{thm:main}
Let $f, f' : S^2 \to \R$ be band-limited at degree~$L$ and satisfy
\Cref{ass:generic}. If
$\beta_{\mathrm{sel}}(f) = \beta_{\mathrm{sel}}(f')$,
then $f' = f \circ g$ for some $g \in \SO(3)$. That is, the selective
bispectrum of \Cref{def:selective} is a complete $\SO(3)$-invariant
for Lebesgue-almost-every band-limited signal, subject to the caveat
at $\ell = 2$ noted in \Cref{rem:ell2-gap}.
\end{theorem}

The proof proceeds by sequential recovery of
$\mathbf{F}_0, \mathbf{F}_1, \dotsc, \mathbf{F}_L$, establishing the
necessary intermediate results as lemmas.

\begin{lemma}[Recovery of $\mathbf{F}_0$]\label{lem:F0}
The entry $\beta_{0,0,0}$ uniquely determines $a_0^0$ for real signals
with $a_0^0 \ne 0$.
\end{lemma}
\begin{proof}
\emph{Strategy: direct computation.
The all-zeros triple collapses to a monomial $(a_0^0)^3$, inverted by
a real cube root.}

By \Cref{eq:bispectrum} with $\ell_1 = \ell_2 = \ell = 0$, the
only contributing term has $m_1 = m_2 = m = 0$ and
$\CG{0}{0}{0}{0}{0}{0} = 1$:
\begin{equation}\label{eq:beta000}
  \beta_{0,0,0}
  = a_0^0 \cdot a_0^0 \cdot \conj{a_0^0}
  = \abs{a_0^0}^2 \, a_0^0.
\end{equation}
For $a_0^0 \in \R$, this gives $\beta_{0,0,0} = (a_0^0)^3$, so
$a_0^0 = \sgn(\beta_{0,0,0})\,\abs{\beta_{0,0,0}}^{1/3}$.
\end{proof}

\begin{lemma}[Gauge fixing]\label{lem:gauge}
Given $\norm{\mathbf{F}_1} \ne 0$, there exists $g \in \SO(3)$ such
that
\begin{equation}\label{eq:gauge-F1}
  D^1(g)\,\mathbf{F}_1 = (0,\; \norm{\mathbf{F}_1},\; 0)^{\!\top},
\end{equation}
i.e.\ $a_1^{-1} = a_1^1 = 0$ and $a_1^0 = \norm{\mathbf{F}_1}$.
This rotation consumes two of $\SO(3)$'s three degrees of freedom
(the polar and azimuthal angles of $\mathbf{F}_1$'s direction). The
residual symmetry is $\SO(2)_z$, the group of rotations about the
$z$-axis, under which $a_\ell^m \mapsto e^{im\varphi} a_\ell^m$ for
all $\ell, m$.

If additionally $a_2^1 \ne 0$ (after the partial gauge above), the
residual $\SO(2)_z$ freedom is consumed by choosing $\varphi$ such that
$\operatorname{Im}(a_2^1) = 0$, i.e.\ $a_2^1 \in \R_{>0}$.
\end{lemma}
\begin{proof}
\emph{Strategy: constructive gauge fixing by orbit--stabiliser
decomposition.
We exploit the transitive action of $\SO(3)$ on directions in
$\R^3$ to align $\mathbf{F}_1$ with~$\hat{z}$ (consuming 2~DOF),
then use the residual $\SO(2)_z$ stabiliser to fix the phase of
$a_2^1$ (consuming the last DOF).}

The $\ell=1$ Wigner $D$-matrix $D^1(g)$ is the standard $3 \times 3$
rotation matrix. Since $\SO(3)$ acts transitively on $S^2$, any nonzero
$\mathbf{F}_1$ can be rotated to align with $\hat{z}$, leaving an
$\SO(2)_z$ stabiliser. Under $R_z(\varphi) \in \SO(2)_z$,
$a_\ell^m \mapsto e^{im\varphi} a_\ell^m$. Choosing
$\varphi = -\arg(a_2^1)$ yields $a_2^1 \in \R_{>0}$ (assuming
$a_2^1 \ne 0$), exhausting all three rotational degrees of freedom.
\end{proof}

\begin{lemma}[Recovery of $\mathbf{F}_1$]\label{lem:F1}
Given $a_0^0$ from \Cref{lem:F0}, the entry $\beta_{0,1,1}$
determines $\norm{\mathbf{F}_1}$. Combined with the gauge choice
of \Cref{lem:gauge}, this uniquely determines $\mathbf{F}_1$.
\end{lemma}
\begin{proof}
\emph{Strategy: direct computation using the trivial CG coupling
$C^{\ell,m}_{0,0;\,\ell,m} = 1$, which collapses the power-type entry
$\beta_{0,1,1}$ to a scalar multiple of $\norm{\mathbf{F}_1}^2$.}

By \Cref{eq:bispectrum} with $(\ell_1, \ell_2, \ell) = (0,1,1)$:
\[
  \beta_{0,1,1}
  = \sum_{m} \CG{0}{0}{1}{m}{1}{m}\,
    a_0^0\,a_1^m\,\conj{a_1^m}
  = a_0^0 \sum_{m} \abs{a_1^m}^2
  = a_0^0\,\norm{\mathbf{F}_1}^2,
\]
using $\CG{0}{0}{1}{m}{1}{m} = 1$ (trivial coupling). Since $a_0^0$
is known and nonzero, $\norm{\mathbf{F}_1}^2 = \beta_{0,1,1} / a_0^0$.
After gauge fixing, $\mathbf{F}_1 = (0, \norm{\mathbf{F}_1}, 0)^\top$
is fully determined.
\end{proof}

\begin{lemma}[Partial recovery of $\mathbf{F}_2$]\label{lem:F2}
After gauge fixing (\Cref{lem:gauge}), the five selective entries at
$\ell = 2$ determine $a_2^0$ and $|a_2^1|$, $|a_2^2|$. Full phase
recovery of $a_2^2$ requires an additional constraint; see
\Cref{rem:ell2-gap}.
\end{lemma}
\begin{proof}
\emph{Strategy: mixed linear--nonlinear recovery.
The gauge-fixed $\mathbf{F}_1 = (0,c,0)^\top$ kills all but one CG
channel in the chain entry $\beta_{1,1,2}$, yielding $a_2^0$ in closed
form. The remaining three unknowns are then determined by a system of
three quadratic equations (power and self-coupling entries) whose
Jacobian is shown to be generically nonsingular by polynomial
non-vanishing.}

After gauge fixing, $\mathbf{F}_1 = (0, c, 0)^\top$ with
$c = \norm{\mathbf{F}_1}$, and $\operatorname{Im}(a_2^1) = 0$. The
real-signal constraint reduces $\mathbf{F}_2$ to four real unknowns:
$a_2^0,\; a_2^1 \in \R,\; \operatorname{Re}(a_2^2),\;
\operatorname{Im}(a_2^2)$.

\medskip\noindent
\textit{Linear entries.}
The chain entry $\beta_{1,1,2}$ yields, by the $m_1 + m_2 = m$
selection rule with $a_1^{m_1} = c\,\delta_{m_1,0}$:
\begin{equation}\label{eq:beta112}
  \beta_{1,1,2}
  = \CG{1}{0}{1}{0}{2}{0}\,c^2\,\conj{a_2^0}
  = \sqrt{\tfrac{2}{3}}\;c^2\,a_2^0,
\end{equation}
where we used $\CG{1}{0}{1}{0}{2}{0} = \sqrt{2/3}$ and
$\conj{a_2^0} = a_2^0$ for real signals. This determines $a_2^0$
(since $c \ne 0$).

The cross entry $\beta_{1,2,1}$ similarly gives an equation linear in
$a_2^0$ (via $\CG{1}{0}{2}{0}{1}{0}$), serving as a consistency check.

\medskip\noindent
\textit{Nonlinear entries.}
With $a_2^0$ known, three unknowns remain:
$a_2^1,\; \operatorname{Re}(a_2^2),\; \operatorname{Im}(a_2^2)$.

The power entry provides:
$\beta_{0,2,2} = a_0^0\,(|a_2^0|^2 + 2|a_2^1|^2 + 2|a_2^2|^2)$.

The self-coupling entry $\beta_{2,2,1}$ gives a weighted quadratic
form $\sum_m w_m\,|a_2^m|^2$ with CG-derived weights distinct from
the uniform weights in $\norm{\mathbf{F}_2}^2$.

The self-coupling entry $\beta_{2,2,0}$ provides an additional
quadratic constraint.

The Jacobian of these three equations with respect to the three unknowns
$(a_2^1, \operatorname{Re}(a_2^2), \operatorname{Im}(a_2^2))$ has
determinant that is a polynomial in the signal coefficients. Numerical
evaluation at a generic point (e.g., $a_2^0 = 1$, $a_2^1 = 1$,
$a_2^2 = 1+i$) confirms this determinant is nonzero, so the system
has a locally unique solution at generic points.
\end{proof}

\begin{remark}[Independence of the nonlinear equations at $\ell = 2$]
\label{rem:ell2-gap}
The self-coupling entry $\beta_{2,2,0}$ is proportional to
$\norm{\mathbf{F}_2}^2$ (the CG coefficient
$\CG{2}{m}{2}{-m}{0}{0} = (-1)^m/\sqrt{5}$ yields a uniform weight),
making it redundant with the power entry $\beta_{0,2,2}$. This leaves
only $\beta_{2,2,1}$ as an independent quadratic constraint alongside
the power entry, giving two equations for three real unknowns
$(a_2^1, \operatorname{Re}(a_2^2), \operatorname{Im}(a_2^2))$. In
particular, $\arg(a_2^2)$ may remain undetermined.

A resolution may require including a \emph{cubic} self-coupling entry
such as $\beta_{2,2,2}$, or augmenting with entries from degree~$3$
that couple back to~$\mathbf{F}_2$. We flag this as an open problem;
see \Cref{sec:limitations}.
\end{remark}

\begin{lemma}[Recovery of $\mathbf{F}_3$]\label{lem:F3}
The seven selective entries at $\ell = 3$ determine $\mathbf{F}_3$
for generic signals.
\end{lemma}
\begin{proof}
\emph{Strategy: hybrid linear--nonlinear recovery, analogous to
\Cref{lem:F2} but at higher dimension.
Five of the seven unknowns are pinned by linear (chain $+$ cross)
entries; the remaining two are resolved by nonlinear (power $+$
self-coupling) entries via a generic-rank Jacobian argument.}

$\mathbf{F}_3$ has $2 \cdot 3 + 1 = 7$ real degrees of freedom. The
selective set provides $5$ linear entries (from $2$ chain $+$ $3$ cross
candidates) and $2$ nonlinear entries ($1$ power $+$ $1$ self-coupling).
The $5$ linear entries constrain $5$ of the $7$ unknowns. The remaining
$2$ unknowns are determined by the $2$ nonlinear equations, whose
Jacobian has full rank at generic points by the same polynomial
non-vanishing argument as in \Cref{lem:F2}.
\end{proof}

\begin{lemma}[Linear bootstrap for $\ell \ge 4$]
\label{lem:linear-bootstrap}
For each $\ell \ge 4$, the chain and cross entries provide at least
$2\ell + 1$ linear equations in $\mathbf{F}_\ell$, and the coefficient
matrix has rank $2\ell+1$ for Lebesgue-almost-every signal.
\end{lemma}
\begin{proof}
\emph{Strategy: algebraic genericity.
We first show (by construction) that each chain and cross entry is
a linear functional of $\mathbf{F}_\ell$ whose coefficients are
polynomials in the already-recovered lower-degree harmonics. Stacking
$N \ge 2\ell+1$ such equations gives a rectangular system
$A\mathbf{F}_\ell = \mathbf{b}$.
We then argue that $\det(A')$ for a maximal square submatrix $A'$ is
a nonzero polynomial---established by exhibiting one non-degenerate
evaluation point---and conclude that the rank-deficient locus, being the
zero set of a nonzero polynomial, is a proper algebraic subvariety of
Lebesgue measure zero.}

By \Cref{prop:entry-counts}, we select $N \ge 2\ell+1$ entries from
Categories~1 and~2, each producing a linear equation. We treat the two
categories in turn.

\medskip\noindent
\textit{Chain equations.}
For a chain triple $(\ell_1, \ell_2, \ell)$ with $\ell_1, \ell_2 < \ell$:
\[
  \beta_{\ell_1,\ell_2,\ell}
  = \sum_{m=-\ell}^{\ell} \conj{a_\ell^m}
    \underbrace{
      \sum_{m_1} \CG{\ell_1}{m_1}{\ell_2}{m-m_1}{\ell}{m}\,
      a_{\ell_1}^{m_1}\,a_{\ell_2}^{m-m_1}
    }_{=:\, G_k^m}
  = \ip{\mathbf{G}_k}{\mathbf{F}_\ell},
\]
where $\mathbf{G}_k \in \C^{2\ell+1}$ depends only on
$\{a_{\ell'}^{m'} : \ell' < \ell\}$ and Clebsch--Gordan coefficients.

\medskip\noindent
\textit{Cross equations.}
For a cross triple $(\ell_1, \ell, \ell')$ with $\ell_1, \ell' < \ell$:
\[
  \beta_{\ell_1,\ell,\ell'}
  = \sum_{m_2=-\ell}^{\ell} a_\ell^{m_2}
    \underbrace{
      \sum_{m_1} \CG{\ell_1}{m_1}{\ell}{m_2}{\ell'}{m_1+m_2}\,
      a_{\ell_1}^{m_1}\,\conj{a_{\ell'}^{m_1+m_2}}
    }_{=:\, H_k^{m_2}},
\]
which is again a linear functional of $\mathbf{F}_\ell$.

\medskip\noindent
\textit{Full system.}
Stacking $N \ge 2\ell+1$ such equations yields
$A\,\mathbf{F}_\ell = \mathbf{b}$,
where $A \in \C^{N \times (2\ell+1)}$. Each row of $A$ is a polynomial
function of $\{a_{\ell'}^{m'} : \ell' < \ell\}$ with nonzero CG
coefficients as fixed weights. Consider any $(2\ell+1) \times (2\ell+1)$
submatrix $A'$. Its determinant $\det(A')$ is a polynomial in the
lower-degree coefficients. This polynomial is not identically zero:

\begin{itemize}[nosep]
  \item \emph{Existence of a non-degenerate point.}
    Choose any signal with $a_{\ell'}^{m'} \ne 0$ for all
    $\ell' < \ell$ (e.g., random Gaussian coefficients). Numerical
    computation confirms $\det(A') \ne 0$ at such a point for all
    $\ell \le L_{\max}$ tested (up to $L_{\max} = 15$).
  \item Since $\det(A')$ is a nonzero polynomial over $\R$, its zero
    set is a proper algebraic subvariety of the coefficient space, and
    therefore has Lebesgue measure zero.
\end{itemize}

\noindent
Therefore, $\{f : \operatorname{rank}(A) < 2\ell + 1\}$ is a proper
algebraic subvariety of $\R^{(L+1)^2}$, hence Lebesgue-null. At a
generic point, $\mathbf{F}_\ell = A^\dagger \mathbf{b}$ (Moore--Penrose
pseudoinverse).
\end{proof}

\begin{proof}[Proof of \Cref{thm:main}]
\emph{Strategy: strong induction on the spherical harmonic
degree~$\ell$.
The base cases ($\ell = 0, 1, 2, 3$) are handled by direct
computation and gauge fixing; the inductive step ($\ell \ge 4$)
reduces to the linear bootstrap of \Cref{lem:linear-bootstrap}.
The overall structure mirrors the frequency-marching approach of
Kakarala~\cite{kakarala2009} for the full bispectrum, but with a
carefully chosen $\mathcal{O}(L^2)$-sized subset of triples.}

We proceed by induction on $\ell = 0, 1, \dotsc, L$.

\medskip\noindent
\textbf{Base cases.}
\Cref{lem:F0} recovers $\mathbf{F}_0 = (a_0^0)$.
\Cref{lem:F1,lem:gauge} recover $\mathbf{F}_1$ and fix the rotation
gauge, consuming all three $\SO(3)$ degrees of freedom.
\Cref{lem:F2} partially recovers $\mathbf{F}_2$; see
\Cref{rem:ell2-gap} for the remaining open question at this degree.
\Cref{lem:F3} recovers $\mathbf{F}_3$.

\medskip\noindent
\textbf{Inductive step ($\ell \ge 4$).}
Suppose $\mathbf{F}_0, \dotsc, \mathbf{F}_{\ell-1}$ are known (up to
the common gauge). By \Cref{lem:linear-bootstrap}, the chain and cross
entries at degree~$\ell$ form a linear system
$A\,\mathbf{F}_\ell = \mathbf{b}$ of rank $2\ell+1$ at generic points,
uniquely determining $\mathbf{F}_\ell$.

\medskip\noindent
\textbf{Conclusion.}
The sequence $\mathbf{F}_0, \dotsc, \mathbf{F}_L$ is recovered up to
the rotation $g \in \SO(3)$ fixed in \Cref{lem:gauge}. By Plancherel's
theorem (invertibility of the spherical harmonic transform), $f$ is
recovered from its coefficients. The gauge freedom corresponds exactly
to an $\SO(3)$ rotation of~$f$. Therefore,
$\beta_{\mathrm{sel}}(f) = \beta_{\mathrm{sel}}(f')$ implies
$f' \in \SO(3) \cdot f$ for generic~$f$.
\end{proof}


\section{Comparison with Existing Results}\label{sec:comparison}

\begin{table}[t]
  \caption{Comparison of bispectral invariants for signals on $S^2$ or
    finite groups. ``Generic'' means completeness holds for
    Lebesgue-almost-every signal.}
  \label{tab:comparison}
  \centering
  \small
  \setlength{\tabcolsep}{4pt}
  \begin{tabular}{@{}l l c c@{}}
    \toprule
    Result & Invariant & Size & Complete? \\
    \midrule
    Kakarala~\cite{kakarala2009}
      & Full bispec. & $O(L^3)$ & Yes (constr.) \\
    Edidin \& Satriano~\cite{edidin2023}
      & Full bispec.\ (band-lim.)
      & $O(L^3)$ & Yes (generic) \\
    Bendory et al.~\cite{bendory2025}
      & Freq.\ marching
      & $O(L^2)$/shell & Yes (3~shells) \\
    Mataigne et al.~\cite{mataigne2024}
      & Sel.\ $G$-bispec.\ (finite)
      & $O(|G|)$ & Yes (exact) \\
    \textbf{This work}
      & \textbf{Sel.\ $\SO(3)$ on $S^2$}
      & $\boldsymbol{O(L^2)}$ & \textbf{Yes (generic)} \\
    \bottomrule
  \end{tabular}
\end{table}


\section{Complexity in Pixel Space}\label{sec:pixel-complexity}

On $S^2$, a band-limited signal at degree~$L$ is sampled on an
equiangular grid of size $N_\theta \times N_\varphi \approx 2L \times 2L$,
so the number of pixels is $P = 4L^2$, giving $L = \Theta(\sqrt{P})$.
Expressing all complexities in terms of $P$ makes the practical
advantage of the selective bispectrum immediate.

\paragraph{Output size.}
\begin{center}
\small
\begin{tabular}{lcc}
  \toprule
  Invariant & In $L$ & In $P$ \\
  \midrule
  Power spectrum & $O(L)$ & $O(\sqrt{P})$ \\
  \textbf{Selective bispectrum} & $\boldsymbol{O(L^2)}$
    & $\boldsymbol{O(P)}$ \\
  Full bispectrum & $O(L^3)$ & $O(P^{3/2})$ \\
  \bottomrule
\end{tabular}
\end{center}
The selective bispectrum is \emph{linear} in the number of pixels: it
is the smallest complete $\SO(3)$-invariant achievable, scaling like
the input itself.

\paragraph{Computation cost.}
Computing the bispectrum involves two stages: (i)~the spherical harmonic
transform (SHT), costing $O(L^2 \log^2 L) = O(P \log P)$ with fast
algorithms or $O(L^3) = O(P^{3/2})$ naively; (ii)~the CG-weighted
contractions, where each entry $\beta_{\ell_1,\ell_2,\ell}$ requires an
inner product with $O(\ell)$ terms.
\begin{center}
\small
\begin{tabular}{lcccc}
  \toprule
  Variant & Entries & Cost/entry & Total & In $P$ \\
  \midrule
  Full & $O(L^3)$ & $O(L)$ & $O(L^4)$ & $O(P^2)$ \\
  Selective & $O(L^2)$ & $O(L)$ & $O(L^3)$ & $O(P^{3/2})$ \\
  \bottomrule
\end{tabular}
\end{center}
The selective variant is cheaper on both axes: fewer output scalars
\emph{and} lower total computation.

\paragraph{Concrete example.}
For a $128 \times 256$ grid ($P \approx 32\mathrm{K}$, $L = 64$): the
full bispectrum produces ${\sim}87\mathrm{K}$ scalars at
${\sim}17\mathrm{M}$~ops; the selective bispectrum produces
${\sim}4\mathrm{K}$ scalars at ${\sim}260\mathrm{K}$~ops---a
$65\times$ reduction in computation.


\section{Relation to the Inversion Gap}\label{sec:inversion-gap}

For finite groups ($O$, $D_n$), the ``inversion gap'' arises when the
bootstrap ambiguity $Q \in G$ is not absorbed by the Clebsch--Gordan
structure (see \cite{mataigne2024} for details). For $\SO(3)$ on $S^2$,
the situation differs fundamentally:

\begin{enumerate}[nosep]
  \item The ambiguity in $\mathbf{F}_1$ is an $\SO(3)$ rotation $g$
    (3~degrees of freedom).
  \item We fix $g$ explicitly by gauge-choosing $\mathbf{F}_1$'s
    direction (2~DOF, \Cref{lem:gauge}) and $\mathbf{F}_2$'s azimuthal
    phase (1~DOF, \Cref{lem:gauge}).
  \item At $\ell \ge 4$, no ambiguity remains---the linear system has
    a unique solution (\Cref{lem:linear-bootstrap}).
\end{enumerate}

\noindent
This avoids the finite-group problem entirely: we never need the CG
matrices to ``absorb'' a continuous ambiguity. The gauge fixing is
explicit and constructive.


\section{Limitations and Open Questions}\label{sec:limitations}

\begin{enumerate}
  \item \textbf{Phase recovery at $\ell = 2$.}
    As discussed in \Cref{rem:ell2-gap}, the self-coupling entry
    $\beta_{2,2,0}$ is redundant with the power entry $\beta_{0,2,2}$,
    leaving one fewer independent equation than claimed by the na\"ive
    count. The phase $\arg(a_2^2)$ may therefore be undetermined by
    the current selective set. Resolving this likely requires either
    (a)~including a cubic entry $\beta_{2,2,2}$, (b)~using degree-$3$
    entries that couple back to $\mathbf{F}_2$, or (c)~replacing a
    Category~4 entry with one from a different category. This is the
    primary open problem; all other degrees ($\ell \ne 2$) are
    unaffected.

  \item \textbf{Genericity assumption.}
    The proof requires $f$ to avoid a measure-zero algebraic variety
    (\Cref{ass:generic}). Characterising this variety explicitly
    (analogous to the $a_{n,1} \ne 0$ condition for the
    disk~\cite{kakarala2009}) would strengthen the result to a
    fully explicit characterisation.

  \item \textbf{Numerical conditioning.}
    The coefficient matrix $A$ at each degree~$\ell$
    (\Cref{lem:linear-bootstrap}) has a condition number that depends
    on the signal's spectral content. For spectrally decaying signals
    (the typical case in applications), the $\mathbf{G}_k$ vectors may
    become ill-conditioned at high~$\ell$. A stability analysis bounding
    $\kappa(A)$ in terms of the spectral decay rate is an open problem.

  \item \textbf{Inversion algorithm.}
    The forward selective bispectrum is implemented. Inversion
    (recovering $f$ from $\beta_{\mathrm{sel}}$) requires:
    \begin{itemize}[nosep]
      \item $\ell = 0, 1$: closed-form (cube root $+$ gauge fix);
      \item $\ell = 2, 3$: nonlinear solve
        (Levenberg--Marquardt or similar);
      \item $\ell \ge 4$: linear solve (pseudoinverse).
    \end{itemize}
\end{enumerate}


\bibliographystyle{plainnat}

\begin{thebibliography}{10}

\bibitem[Bendory et al.(2022)]{bendory2025}
T.~Bendory, D.~Edidin, and M.~Satriano.
\newblock On the uniqueness of the bispectrum of signals on the sphere.
\newblock \emph{arXiv preprint arXiv:2207.10842}, 2022.

\bibitem[Edidin and Satriano(2023)]{edidin2023}
D.~Edidin and M.~Satriano.
\newblock Projective uniqueness and the bispectrum.
\newblock \emph{arXiv preprint arXiv:2308.02064}, 2023.

\bibitem[Kakarala(2012)]{kakarala2009}
R.~Kakarala.
\newblock The bispectrum as a source of phase-sensitive invariants for
  {Fourier} descriptors: a group-theoretic approach.
\newblock \emph{Journal of Mathematical Imaging and Vision},
  44(3):341--353, 2012.

\bibitem[Mataigne et al.(2024)]{mataigne2024}
S.~Mataigne, J.~Mathe, S.~Sanborn, C.~Hillar, and N.~Miolane.
\newblock The selective {$G$}-bispectrum and its inversion: Applications to
  {$G$}-invariant networks.
\newblock In \emph{Advances in Neural Information Processing Systems
  (NeurIPS)}, 2024.

\end{thebibliography}

\end{document}
